\section{The Topic: Accountability for Heritage Destruction in Armed Conflict}

\subsection{Key Terms}

\textbf{Cultural heritage} = the physical and intangible legacy that communities inherit from past generations and pass on to future ones. This includes tangible heritage and intangible heritage. UNESCO's legal framework covers both, though the two are protected by different conventions and enjoy very different levels of enforcement.

\textbf{Tangible cultural heritage} = includes monuments, archaeological sites, museums, historic buildings, libraries, archives. Can be further divided into movable and immovable heritage.

\textbf{Intangible cultural heritage} = heritage forming living traditions. Includes oral histories, languages, music, rituals, and craft practices.

\textbf{World Heritage Sites} = properties inscribed on UNESCO's World Heritage List under the 1972 Convention as possessing "Outstanding Universal Value". Their significance transcends borders and belongs to all of humanity. Inscription does not guarantee protection; it does create legal obligations for States of protection.

\textbf{1945 Hague Convention} = cornerstone international treaty for heritage protection in armed conflict. It obliges States to safeguard cultural property in times of war and prohibits deliberate attacks on monuments, museums, and historic sites. The Convention introduced the Blue Shield emblem, the cultural equivalent of the Red Cross  to mark protected sites.\footnote{Blue Shield International. (2020). \emph{About Blue Shield}.\url{https://theblueshield.org/about-us/}}

\textbf{International Humanitarian Law (IHL)} = body of rules that seek to limit the effects of armed conflict. It protects civilians, combatants who are no longer fighting, and cultural property. Both the Hague Regulations and the Geneva Conventions include provisions on cultural heritage.\footnote{Henckaerts, J. M., \& Doswald-Beck, L. (2005). \emph{Customary international humanitarian law, Volume I: Rules}. Cambridge University Press / ICRC.\url{https://doi.org/10.1017/CBO9780511804700}}

\textbf{War crime} = serious violation of the law and customs applicable in armed conflict, defined under the Rome Statute of the International Criminal Court (ICC).\footnote{International Criminal Court. (1998). \emph{Rome Statute of the International Criminal Court}.\url{https://www.icc-cpi.int/resource-library/documents/rs-eng.pdf}} Since the 2016 conviction of Ahmad Al Faqi Al Mahdi, the intentional direction of attacks against protected cultural or religious buildings can constitute a war crime.\footnote{International Criminal Court. (2016). \emph{Prosecutor v. Ahmad Al Faqi Al Mahdi, Judgment and Sentence} (ICC-01/12-01/15).\url{https://www.icc-cpi.int/mali/al-mahdi}} However, this landmark ruling only established individual criminal responsibility. It left the question of State responsibility largely unanswered.\footnote{Wierczyńska, K., \& Jakubowski, A. (2017). Individual responsibility for deliberate destruction of cultural heritage: Contextualising the ICC judgment in the Al-Mahdi case. \emph{Chinese Journal of International Law, 16}(4), 695–721.\url{https://doi.org/10.1093/chinesejil/jmx029}} \footnote{Schabas, W. A. (2016). \emph{The International Criminal Court: A commentary on the Rome Statute} (2nd ed.). Oxford University Press.}\par \textbf{State responsibility} = States can be held accountable for internationally wrongful acts attributable to it including the destruction of cultural heritage carried out by its armed forces.   This is distinct from criminal responsibility, which attaches to individuals.\footnote{International Law Commission. (2001). \emph{Articles on Responsibility of States for Internationally Wrongful Acts}. United Nations.\url{https://legal.un.org/ilc/texts/instruments/english/draft_articles/9_6_2001.pdf}}

\textbf{Non-State armed groups (NSAG)} = armed organizations that operate outside the direct control of a recognized state. These organisations carried out some of the most systematic heritage destruction in recent history, yet they fall outside the traditional framework of state responsibility and are notoriously difficult to bring before international courts.\footnote{Lostal, M. (2017). \emph{International cultural heritage law in armed conflict: Case studies of Syria, Libya, Mali, the invasion of Iraq, and the Buddhas of Bamiyan}. Cambridge University Press.}

\textbf{Iconoclasm} = the destruction of images refers to the deliberate elimination of cultural or religious symbols, often as an assertion of ideological power. It is a targeted, intentional act designed to erase the identity, memory, and legitimacy of a people.\footnote{Francioni, F., \& Lenzerini, F. (2003). The destruction of the Buddhas of Bamiyan and international law. \emph{European Journal of International Law, 14}(4), 619–651.\url{https://doi.org/10.1093/ejil/14.4.619}}

\textbf{Memorialisation} = the practices of reconstruction, documentation, commemoration, oral testimony through which communities process and preserve the memory of what was lost.

\footnote{Isakhan, B., \& Meskell, L. (2019). UNESCO's project to 'Revive the Spirit of Mosul': Iraqi and Syrian opinion on heritage reconstruction after the Islamic State. \emph{International Journal of Heritage Studies, 25}(11), 1189–1204.\url{https://doi.org/10.1080/13527258.2019.1578988}} \footnote{UNESCO World Heritage Centre. (n.d.). \emph{Reconstruction of the destroyed mausoleums of Timbuktu (Mali)}.\url{https://whc.unesco.org/en/canopy/timbuktu/}}

\subsection{Introduction to the Topic}

The central, painful truth that this committee must confront is this: the international community has known for decades that heritage destruction is wrong, has written laws against it, has condemned it in resolutions and speeches and has almost entirely failed to stop it or meaningfully punish those responsible.\footnote{Lostal, M. (2017). \emph{International cultural heritage law in armed conflict: Case studies of Syria, Libya, Mali, the invasion of Iraq, and the Buddhas of Bamiyan}. Cambridge University Press.}

Cultural heritage is not simply a collection of old objects and pretty buildings. It is the physical embodiment of identity. When the Buddhas of Bamiyan were dynamited by the Taliban in 2001, Afghanistan did not just lose two statues, it lost irreplaceable proof of a pre-Islamic, cosmopolitan past that challenged the Taliban's entire ideology of cultural purity.\footnote{Francioni, F., \& Lenzerini, F. (2003). The destruction of the Buddhas of Bamiyan and international law. \emph{European Journal of International Law, 14}(4), 619–651.\url{https://doi.org/10.1093/ejil/14.4.619}} When ISIS bulldozed the ancient city of Nimrud in 2015, they did not merely destroy Assyrian ruins, they deliberately severed living people from the oldest layers of their own history.\footnote{UNESCO. (2015b). \emph{UNESCO Director-General condemns destruction of Nimrud in Iraq}.\url{https://www.unesco.org/en/articles/unesco-director-general-condemns-destruction-nimrud-iraq}} \footnote{ASOR Cultural Heritage Initiatives. (2015). \emph{Cultural heritage destruction and looting in Syria and Iraq: Monthly report}. American Schools of Oriental Research.\url{https://www.asor.org/chi/}}

Heritage destruction is almost never purely incidental. It is frequently a weapon used to demoralize populations, delegitimize national identities, erase evidence of prior occupancy, and signal total domination. Understanding it as such changes everything about how accountability should be designed.\footnote{UN Special Rapporteur on Cultural Rights. (2016). \emph{Report on the intentional destruction of cultural heritage} (A/71/317). United Nations.\url{https://undocs.org/A/71/317}}

The international legal framework that exists to prevent this is substantial, in theory. The 1954 Hague Convention was the first comprehensive treaty dedicated to cultural property protection in armed conflict  born directly from the ashes of a war in which Nazi Germany systematically looted and destroyed the cultural wealth of occupied Europe.\footnote{Boylan, P. J. (1993). \emph{Review of the Convention for the Protection of Cultural Property in the Event of Armed Conflict (The Hague Convention of 1954)}. UNESCO.\url{https://unesdoc.unesco.org/ark:/48223/pf0000102816}} Its two additional Protocols extended protections and tightened obligations.\footnote{UNESCO. (1999). \emph{Second Protocol to the Hague Convention of 1954 for the Protection of Cultural Property in the Event of Armed Conflict}.\url{https://en.unesco.org/protecting-heritage/convention-and-protocols/second-protocol}} The Rome Statute of the International Criminal Court criminalizes attacks on cultural property as a war crime. UN Security Council resolutions, most notably Resolution 2347 of 2017, the first ever adopted exclusively on the protection of cultural heritage have explicitly named heritage destruction as a threat to international peace and security.\footnote{United Nations Security Council. (2017). \emph{Resolution 2347} (S/RES/2347).\url{https://undocs.org/S/RES/2347(2017)}} The gap between the law on paper and the law in practice is one of the most striking failures of the post-World War II international order.\footnote{Lostal, M. (2017). \emph{International cultural heritage law in armed conflict: Case studies of Syria, Libya, Mali, the invasion of Iraq, and the Buddhas of Bamiyan}. Cambridge University Press.}

Enforcement of international humanitarian law is extraordinarily difficult during active conflict.\footnote{International Committee of the Red Cross. (2005). \emph{Customary international humanitarian law, Volume I: Rules}. Cambridge University Press.\url{https://www.icrc.org/en/doc/assets/files/other/customary-international-humanitarian-law-i-icrc-eng.pdf}} The very moment when heritage is most at risk, is the moment when international monitors have the least access. Accountability mechanisms designed for interstate conflicts are poorly adapted to non-state armed groups, who command no territory recognized by international law and whose members may never appear before any court. Political will is uneven and selective. The destruction of heritage in the Global South has historically received less international attention than equivalent losses in Europe.\footnote{Maurel, C. (2017, January 11). The unintended consequences of UNESCO world heritage listing. \emph{The Conversation}.\url{https://theconversation.com/the-unintended-consequences-of-unesco-world-heritage-listing-71047}} The fundamental architecture of the international system still treats sovereignty as a near-sacred principle, making it extraordinarily difficult to compel States to answer for what their armed forces do.\footnote{International Law Commission. (2001). \emph{Articles on Responsibility of States for Internationally Wrongful Acts}. United Nations.\url{https://legal.un.org/ilc/texts/instruments/english/draft_articles/9_6_2001.pdf}}

\subsection{Historical Evolution of Heritage Protection in Armed Conflict}

The destruction of cultural heritage during armed conflict is not a modern phenomenon. However, what has changed over time is the degree to which the international community has been willing to recognise it as a legal problem, which requires a collective response.

The earliest attempts to regulate wartime conduct date to the 19th Century. The Lieber Code of 1863 prohibited deliberate destruction of property without military necessity and included early protections for museums and libraries.\footnote{Lieber, F. (1863). \emph{Instructions for the Government of Armies of the United States in the Field} (General Orders No. 100). United States War Department.} The Hague Regulations of 1899 and 1907 subsequently established, in a binding multilateral form, that attacks on historic monuments and institutions dedicated to the arts and sciences were forbidden.\footnote{Hague Convention IV Respecting the Laws and Customs of War on Land. (1907). Articles 27 and 56.\url{https://ihl-databases.icrc.org/en/ihl-treaties/hague-regs-1907}} However, these instruments lacked enforcement mechanisms and applied only to interstate conflicts. Their limitations were exposed during the First World War, when the shelling of Reims Cathedral and the burning of Leuven's medieval library demonstrated that even emblematic European heritage enjoyed no effective protection in practice.\footnote{Boylan, P. J. (1993). \emph{Review of the Convention for the Protection of Cultural Property in the Event of Armed Conflict (The Hague Convention of 1954)}. UNESCO.\url{https://unesdoc.unesco.org/ark:/48223/pf0000102816}}

The systematic looting and destruction carried out by the National Socialist regime across occupied Europe exposed the critical limitations of existing international frameworks, providing the direct impetus for the 1954 Hague Convention for the Protection of Cultural Property in the Event of Armed Conflict, the foundational instrument in this field.\footnote{UNESCO. (1954). \emph{Convention for the Protection of Cultural Property in the Event of Armed Conflict}.\url{https://en.unesco.org/protecting-heritage/convention-and-protocols/1954-convention}} For the first time, an international treaty placed heritage protection within a dedicated legal framework linked to a specific intergovernmental organisation.

The 1954 Convention was later supplemented by a Second Protocol in 1999, which introduced enhanced protection for the most significant sites and, critically, established individual criminal responsibility for serious violations, a development whose importance would only become fully apparent with the rise of non-state actors in subsequent decades.\footnote{UNESCO. (1999). \emph{Second Protocol to the Hague Convention of 1954 for the Protection of Cultural Property in the Event of Armed Conflict}.\url{https://en.unesco.org/protecting-heritage/convention-and-protocols/second-protocol}} The same period saw the adoption of the 1972 World Heritage Convention, which reframed heritage protection as a matter of shared human interest rather than purely national concern, and the 1970 Convention addressing illicit trafficking in cultural objects, which frequently accompanies and finances armed conflict.\footnote{UNESCO. (1972). \emph{Convention Concerning the Protection of the World Cultural and Natural Heritage}.\url{https://whc.unesco.org/en/convention/}} \footnote{UNESCO. (1970). \emph{Convention on the Means of Prohibiting and Preventing the Illicit Import, Export and Transfer of Ownership of Cultural Property}.\url{https://www.unesco.org/en/legal-affairs/convention-means-prohibiting-and-preventing-illicit-import-export-and-transfer-ownership}}

The conflicts of the 1990s exposed the inadequacy of this architecture with renewed force. The deliberate shelling of Dubrovnik's Old Town in 1991 and the dynamiting of the Stari Most bridge in Mostar in 1993 demonstrated that heritage destruction was being deployed as a deliberate instrument of ethnic cleansing, designed to erase the physical evidence of coexistence and minority presence.\footnote{UNESCO. (2004). \emph{Stari Most — Mostar's Old Bridge}.\url{https://whc.unesco.org/en/list/946}} The International Criminal Tribunal for the former Yugoslavia subsequently prosecuted individuals for attacks on cultural property, establishing a precedent for individual criminal accountability that directly informed the Rome Statute. \footnote{International Criminal Tribunal for the Former Yugoslavia. (2001). \emph{Prosecutor v. Blaškić} (Case No. IT-95-14).\url{https://www.icty.org/en/case/blaskic}} Adopted in 1998, the Rome Statute of the ICC explicitly criminalised intentional attacks on historic monuments and religious and cultural buildings as war crimes in both international and non-international armed conflicts, addressing the critical gap left by the 1954 Convention, which had been designed primarily for interstate conflicts.\footnote{International Criminal Court. (1998). \emph{Rome Statute of the International Criminal Court}.\url{https://www.icc-cpi.int/resource-library/documents/rs-eng.pdf}}

The destruction of the Bamiyan Buddhas by the Taliban in 2001 exposed yet another structural gap: no mechanism existed to respond to heritage destruction carried out by a governing authority on its own territory outside of an armed conflict with another state.\footnote{Francioni, F., \& Lenzerini, F. (2003). The destruction of the Buddhas of Bamiyan and international law. \emph{European Journal of International Law, 14}(4), 619–651.\url{https://doi.org/10.1093/ejil/14.4.619}} This episode accelerated a normative shift in which heritage destruction came to be understood as a threat to international peace and security, a framing that ultimately produced UN Security Council Resolution 2347 in 2017, the first resolution dedicated exclusively to cultural heritage protection, formally recognising intentional destruction as a tactic of war.\footnote{United Nations Security Council. (2017). \emph{Resolution 2347} (S/RES/2347).\url{https://undocs.org/S/RES/2347(2017)}}
